\documentclass[11pt]{article}
\usepackage{amsmath,amssymb,amsthm,hyperref}
\newtheorem{theorem}{Theorem}
\newtheorem{lemma}{Lemma}
\newtheorem{corollary}{Corollary}
\newcommand{\E}{\mathbb E}
\newcommand{\Var}{\operatorname{Var}}
\title{Quantitative diagonalization of monotone comparison tensors}
\author{Asymptotic research note; independently reviewed}
\date{30 September 2026}
\begin{document}\maketitle
Let $p_1,\ldots,p_m$ be nondecreasing $[0,1]$-valued functions on a
probability space with an ordered finite support. Suppose
\[
 \E\prod_{j\in S}p_j=\frac1{|S|+1}\qquad(S\subseteq[m]).
\]
Write $\rho=1/12$ and $d_j=\Var(p_j)-\rho$.
These statements hold for arbitrary positive row weights, except where
an equal-weight quadrature interpretation is explicitly mentioned.

\begin{theorem}[Variance defects]
At most one $d_j$ is negative. If the positive defects in decreasing
order are $d_{(1)},d_{(2)},\ldots$, then
\[
 d_{(r)}\le\frac{r}{4L(r-L+1)}\le\frac1r,
 \qquad L=\left\lfloor\frac{r+1}{2}\right\rfloor.
\]
Consequently $\sum_j d_j\le 1+\log m$ and, for
$\bar p=m^{-1}\sum_j p_j$,
\[
 \frac1m\sum_j\E(p_j-\bar p)^2
 =\frac{m-1}{m^2}\sum_j d_j
 \le\frac{1+\log m}{m}.
\]
\end{theorem}
\begin{proof}
The centered Gram matrix is $G=D+\rho\mathbf1\mathbf1^T$.
Two negative defects would give a negative quadratic form on the
vector supported on their two positions with entries $1,-1$.

Represent the ordered probability space on $[0,1]$ by assigning each
atom a subinterval of length equal to its weight and extending each
$p_j$ constantly. Project these functions onto functions constant on
$L$ equal-length subintervals. For one column, let $a_\ell$ be its
oscillation on subinterval $\ell$, ignoring endpoints of measure zero.
Monotonicity gives $\sum_\ell a_\ell\le1$. Its squared projection
error is at most
\[
 \frac1{4L}\sum_\ell a_\ell^2\le\frac1{4L},
\]
by the elementary variance bound for a variable in an interval.

Choose $r$ columns with defects at least $\delta=d_{(r)}>0$.
Their Gram matrix is at least $\delta I_r$. The projected centered
columns span a space of dimension at most $L-1$. By orthogonal
projection, or the variational principle for singular values, their
total squared projection error is at least $(r-L+1)\delta$.
Its upper bound is $r/(4L)$, proving the estimate. The harmonic-sum
bound follows. Finally direct expansion of the centered Gram matrix
gives the displayed identity for the average distance to $\bar p$.
\end{proof}

\begin{lemma}[Elementary symmetric means]\label{sampling}
For $x_1,\ldots,x_r\in[0,1]$, put
$t=r^{-1}\sum_i x_i$ and $V=r^{-1}\sum_i(x_i-t)^2$.
For $2\le s\le r$,
\[
 0\le t^s-\frac{e_s(x_1,\ldots,x_r)}{\binom rs}
 \le\frac{\binom s2}{r-1}V.
\]
\end{lemma}
\begin{proof}
Replace any pair $x,y$ by two copies of $(x+y)/2$. This increases
$e_s$ by $((x-y)^2/4)e_{s-2}$ of the remaining $r-2$ coordinates,
and decreases $V$ by $(x-y)^2/(2r)$. Since all remaining coordinates
are in $[0,1]$,
\[
 0\le\Delta\frac{e_s}{\binom rs}
 \le\frac{\binom s2}{r-1}(-\Delta V).
\]
Repeated pair averaging can be chosen to converge to the constant
vector $(t,\ldots,t)$: for example, always average a smallest and a
largest entry; unless $V=0$, its decrease is at least $V/(2r)$.
Telescoping and continuity prove both inequalities.
\end{proof}

\begin{theorem}[A scalar approximate rule from many columns]\label{approx}
Let $I$ be any set of $r\ge2$ columns and put
$t=r^{-1}\sum_{j\in I}p_j$ and $\Delta=\sum_{j\in I}d_j$.
Then $\Delta\ge0$, $\E t=1/2$, and simultaneously for $2\le s\le r$,
\[
 0\le\E t^s-\frac1{s+1}
 \le\binom s2\frac{\Delta}{r^2}.
\]
In particular, remove $R$ columns with largest defects, where
$1\le R\le m-2$, and average the remaining $r=m-R$ columns.
Then
\[
 0\le\E t^s-\frac1{s+1}
 \le\frac{\binom s2}{Rr}\qquad(2\le s\le r).
\]
For $R=\lfloor m/2\rfloor$, this is $O(s^2/m^2)$ through a degree
of order $m$. If the original rows have equal weights, the values of
$t$ form an equal-weight scalar approximate quadrature rule on the
same number of nodes.
\end{theorem}
\begin{proof}
At each row apply Lemma~\ref{sampling} to the selected columns.
The expectation of their normalized elementary symmetric polynomial
is exactly $1/(s+1)$. Their mean squared dispersion is
\[
 \E V=\frac1r\sum_{j\in I}\E p_j^2-\E t^2
 =\frac{r-1}{r^2}\Delta,
\]
so $\Delta\ge0$ and the first estimate follows. Every remaining
positive defect after deletion is at most $1/R$ by the preceding
theorem; negative defects only decrease $\Delta$. Hence
$\Delta\le r/R$, yielding the second estimate.
\end{proof}


\begin{lemma}[A projection adapted to a small mean]\label{smallmean}
Let $q$ be a nonincreasing $[0,1]$-valued function on $[0,1]$ with
$\int q=\alpha>0$. There is a partition into $L$ intervals, depending
only on $\alpha$ and $L$, such that its orthogonal step projection $Pq$
satisfies
\[
 \|q-Pq\|_2^2\le\frac{4\alpha}{L(1+\alpha)}.
\]
The same partition works for every such $q$.
\end{lemma}
\begin{proof}
Partition $u=t/(\alpha+t)$ into $L$ equal increments from $0$ to
$1/(1+\alpha)$, and pull the partition back to $[0,1]$.
For a step $h_s(t)=\mathbf1_{[0,s]}(t)$ with $s$ in a partition
interval $[a,b]$, its squared projection error is
\[
 \|h_s-Ph_s\|_2^2
 =\frac{(s-a)(b-s)}{b-a}\le s-a
 \le\frac{(\alpha+s)^2}{\alpha L(1+\alpha)}.
\]
The last inequality follows from
$s-a=(\alpha+s)(\alpha+a)(u(s)-u(a))/\alpha$.
Write $q=q(1)+\int h_s\,d\nu(s)$ with $\nu$ a positive measure.
Then $\nu([0,1])\le1$ and $\int s\,d\nu(s)\le\alpha$.
The triangle inequality in $L^2$ gives
\[
 \|q-Pq\|_2
 \le\frac{\int(\alpha+s)\,d\nu(s)}
 {\sqrt{\alpha L(1+\alpha)}}
 \le2\sqrt{\frac{\alpha}{L(1+\alpha)}}.
\]
This proof also covers finite step functions and atoms exactly.
\end{proof}

\begin{corollary}[Endpoint-tilted diagonalization]\label{tilted}
Fix a set $S$ of $k$ columns and reweight the rows by
\[
 d\mathbb P_S=(k+1)\prod_{i\in S}p_i\,d\mathbb P.
\]
The remaining columns have mixed moments
$\mu_s^{(k)}=(k+1)/(k+s+1)$. Define their variance defects by
\[
 d_j^{(k)}=\E_S p_j^2-\frac{k+1}{k+3}.
\]
The $r$th largest positive defect satisfies
\[
 d_{(r)}^{(k)}\le\frac{16}{(k+2)r}.
\]
After deleting $R\ge1$ remaining columns with largest defects and
averaging the other $r=m-k-R\ge2$ columns to obtain $t$, one has,
simultaneously for $2\le s\le r$,
\[
 0\le\E_S t^s-\frac{k+1}{k+s+1}
 \le\frac{16\binom s2}{(k+2)Rr}.
\]
In particular, when $k,R,r$ are all of order $m$, the absolute
monomial errors are $O(s^2/m^3)$ through degree of order $m$.
\end{corollary}
\begin{proof}
Rows assigned zero new weight can simply be omitted. Mixed moments
follow by adjoining $S$ to any selected set of remaining columns.
The common mean is $(k+1)/(k+2)$ and the off-diagonal covariance is
\[
 \rho_k=\frac{k+1}{(k+2)^2(k+3)}.
\]
Apply Lemma~\ref{smallmean} to $q_j=1-p_j$, of common mean
$\alpha=1/(k+2)$, using the ordered probability coordinate for the
new row weights. A subfamily of $u$ columns whose defects are at least
$\delta>0$ has Gram matrix at least $\delta I_u$.
The projection argument therefore gives
\[
 (u-L+1)\delta\le\frac{4u\alpha}{L(1+\alpha)}.
\]
Choose $L=\lfloor(u+1)/2\rfloor$ and use
$L(u-L+1)\ge u^2/4$ to obtain $\delta\le16\alpha/u$.
Finally Lemma~\ref{sampling} and the exact dispersion identity from
Theorem~\ref{approx} apply unchanged to the beta moments. The sum of
the retained defects is at most $16r/((k+2)R)$, proving the claim.
\end{proof}

For example, any two columns left after the deletion in Theorem~\ref{approx} satisfy
$\E(p_i-p_j)^2=d_i+d_j\le2/R$. None of these conclusions assumes
strict increase or rules out flat columns. The moment error in
Theorem~\ref{approx} is an absolute error for monomials. It does not
supply a small error for arbitrary degree-$r$ polynomials, whose
monomial coefficients can be large. In particular this note does not
claim the universal quadratic lower bound on the number of rows, or
an exact scalar quadrature replacement.
\end{document}
